\section{Introdução}

A implementação de uma arquitetura de computadores em linguagem de descrição de hardware (HDL) constitui uma etapa fundamental para a consolidação dos conceitos de organização e projeto lógico digital \cite{patterson_mips}. O conjunto de instruções MIPS caracteriza-se por sua regularidade ortogonal e simplicidade de decodificação, tornando-o um modelo de referência amplamente adotado tanto no meio acadêmico quanto em projetos de sistemas embarcados.

Nesta segunda etapa do projeto, o objetivo central consiste na especificação, implementação estrutural e validação funcional do caminho de dados e da unidade de controle do processador MIPS em sua versão monociclo. Em relação à primeira entrega --- focada na especificação do software base e no mapeamento do algoritmo Bubble Sort ---, esta fase consolida os blocos de hardware necessários para executar o fluxo completo de instruções, incluindo desvios condicionais e saltos incondicionais (\textit{jumps}).
